\appendix

\section{附录：符号说明}
为便于阅读，表~\ref{tab:symbols} 汇总本文中频繁使用的主要符号及其含义。

\begin{longtable}{p{0.22\linewidth} p{0.72\linewidth}}
\toprule
符号 & 含义 \\
\midrule
$x_i$ & 第 $i$ 个输入临床文本样本。 \\
$\mathbf{y}_i$ & 第 $i$ 个样本对应的多标签目标向量。 \\
$E(\cdot)$ & 长上下文编码器。 \\
$\mathbf{H}_i$ & 编码器输出的 token 级隐藏状态矩阵。 \\
$\mathbf{z}_i$ & 经池化得到的样本级特征表示。 \\
$\boldsymbol{\mu}_i$ & 输出层预测的标签均值向量。 \\
$\boldsymbol{\sigma}_i^2$ & 输出层预测的输入相关方差向量。 \\
$\mathbf{w}_i$ & 由方差诱导得到的样本权重向量。 \\
$\beta$ & Beta-NLL 中控制稳定化强度的超参数。 \\
$\lambda$ & WLS 精修中的正则系数。 \\
$\mathbf{H}$ & 最后一层前的隐藏表示矩阵。 \\
$\mathbf{W}_j$ & 第 $j$ 个输出维度对应的对角权重矩阵。 \\
$\boldsymbol{\beta}_j$ & 第 $j$ 个输出维度的 WLS 参数向量。 \\
\bottomrule
\caption{本文主要符号说明。}
\label{tab:symbols}
\end{longtable}

\section{附录：两阶段训练流程伪代码}
\begin{algorithm}[H]
\caption{\method{} 的两阶段训练流程}
\KwIn{冻结编码器 $E$，训练集 $\mathcal{D}_{train}$，验证集 $\mathcal{D}_{val}$，超参数 $\beta$ 与 $\lambda$}
\KwOut{训练后的可信输出层参数}
使用编码器 $E$ 对训练集、验证集和测试集提取特征表示；
初始化均值分支 $g_{\mu}$ 与方差分支 $g_{\sigma}$；
\For{epoch $=1,2,\dots,T$}{
  在小批量上前向计算预测均值与对数方差；
  计算 Beta-NLL 损失并更新模型参数；
  在验证集上监控 weighted F1 与 MSE；
  若指标提升，则保存当前模型状态；
}
载入验证集最优参数；
在整个训练集上计算样本权重；
提取均值分支最后一层前的隐藏表示矩阵；
\For{每个输出维度 $j$}{
  求解带正则项的加权最小二乘问题，得到 $\hat{\boldsymbol{\beta}}_j$；
}
用求解得到的参数更新最后线性层；
在测试集上输出分类指标、误差指标和预测方差统计量。
\end{algorithm}

\section{附录：补充实验建议}
为进一步完善本文结论，后续实验可从以下三个方面展开。第一，进行多随机种子重复实验，报告均值和标准差，以验证主结果的稳定性。第二，围绕 pooling 策略、隐藏层维度、Beta-NLL 参数和 WLS 正则系数开展系统超参数分析，以观察各模块对性能和方差分布的影响。第三，增加样本分层分析，例如按预测方差分位点比较错误率、F1 和 AUROC，从而更严格地验证方差输出与样本难度之间的关系。

\section{附录：案例分析展示模板}
正式定稿时，可在附录中加入典型样本分析，以补充正文中的统计结果。建议每个案例至少包含以下信息：输入文本片段、真实标签、模型预测标签、预测分数、预测方差，以及对成功或失败原因的简要解释。按照这一模板，可分别选取低方差正确样本、高方差正确样本和高置信错误样本各一到两个，从而更完整地呈现模型行为边界。
